\section{The Age Profile of Intergenerational Persistence}\label{sec:ageprofile}

Figure~\ref{fig:headline} summarises the section. Panel A traces the log IGE of sons' own earnings at each age from 25
to 50 by race, from one regression that interacts parental income with a quartic in the son's age
\citep{lee2009trends}: the two curves cross in the late twenties---the White--Black gap is $-0.031$ (0.045) at 27 and
$-0.002$ (0.044) at 28---and separate through the thirties, to 0.139$^{***}$ at 35, 0.232$^{***}$ at 40 and
0.317$^{***}$ at 45; between 30 and 45 the gap rises by 0.271$^{***}$ (0.066). Panels B and C place the estimates of
\citet{chetty2020race} on the rank profile: their 31--37 rank--rank slopes (0.28 for Black and 0.32 for White children)
and expected ranks at the 25th parental percentile (32.6 and 45.2) are one point on profiles that are flat in age. The
PSID expected ranks lie a few points below theirs for both races, because the sample and income concept differ, but the
gap is the same (12.9 against 12.6).\footnote{Panel A: \texttt{output/ige\_by\_age.md}, Section 3; panels B and C:
Section 4a (both sexes). Code: \texttt{26.IGEByAge.py}; figure: \texttt{31.PaperBFigure.py}.}

\begin{figure}[!htbp]
\centering
\includegraphics[width=\textwidth]{paperB_age_profile.pdf}
\caption{The age profile of intergenerational persistence, with \citet{chetty2020race} placed on it (PSID)}
\label{fig:headline}
\begin{minipage}{\textwidth}\footnotesize \textit{Notes:} Panel A: log IGE of sons' own earnings at each age, Black
and White, from one Lee--Solon regression with parental income interacted with a quartic in the son's age (child-years
at 25--55; ages above 50 not plotted), 95\% confidence bands. Panel B: rank--rank slope of household income by age band,
sons and daughters, weighted. Panel C: expected household income rank of children whose parents are at the 25th
percentile, by age band. Diamonds: \citet{chetty2020race} at 31--37. Shaded: the 31--37 window. Code:
\texttt{26.IGEByAge.py} (estimates) and \texttt{31.PaperBFigure.py} (figure).\end{minipage}
\end{figure}

\subsection{Persistence rises steeply until the mid-thirties}

Pooled across races and weighted, the IGE of sons' own earnings is 0.276$^{***}$ (0.024) at 25--29, 0.403$^{***}$ at
30--34, 0.454$^{***}$ at 35--39 and 0.492$^{***}$ at 40--44, and that of their family income goes from 0.301$^{***}$
(0.023) to 0.530$^{***}$ (Table~\ref{tab:ige_sons}). The rise is concentrated before the mid-thirties, and it is the
intergenerational face of the life-cycle bias of \citet{haider2006life}: following the same sons from 25--29 to 35--39,
their family income grows faster with parental income by 0.210$^{***}$ (0.031), whereas between 30--34 and 40--44 the
growth gradient is 0.021 (0.036).\footnote{Appendix Table~\ref{tab:growth}. Code: \texttt{26.IGEByAge.py}.} The same
rise appears in the NLSY79 (trend 0.092$^{***}$ per ten years) and the NLSY97 (0.063$^{***}$ on 25--44), and for
daughters' family income (Appendix Tables~\ref{tab:ige_daughters} and \ref{tab:nlsyage}).

\begin{table}[!htbp]
\centering
\caption{Log IGE of sons' income by age at measurement (PSID)}
\label{tab:ige_sons}
\resizebox{\ifdim\width>\textwidth\textwidth\else\width\fi}{!}{%
\resizebox*{!}{\ifdim\totalheight>0.85\textheight 0.85\textheight\else\totalheight\fi}{%
\begin{tabular}{lcc}
\toprule
Age band / group & \begin{tabular}[b]{@{}c@{}}ln own earnings\\(child-years)\end{tabular} & \begin{tabular}[b]{@{}c@{}}ln family income\\(child-years)\end{tabular} \\
 & (1) & (2) \\
\midrule
\multicolumn{3}{l}{\textit{Age 25--29}} \\
Pooled (weighted) & $0.276^{***}$ & $0.301^{***}$ \\
 & (0.024) & (0.023) \\
Black & $0.262^{***}$ & $0.266^{***}$ \\
 & (0.036) & (0.034) \\
White & $0.221^{***}$ & $0.250^{***}$ \\
 & (0.028) & (0.026) \\
White $-$ Black & $-0.041$ & $-0.016$ \\
 & (0.045) & (0.043) \\
Children (Black / White) & 1,607 / 2,474 & 1,662 / 2,506 \\
\addlinespace
\multicolumn{3}{l}{\textit{Age 30--34}} \\
Pooled (weighted) & $0.403^{***}$ & $0.451^{***}$ \\
 & (0.027) & (0.027) \\
Black & $0.282^{***}$ & $0.283^{***}$ \\
 & (0.037) & (0.040) \\
White & $0.369^{***}$ & $0.419^{***}$ \\
 & (0.031) & (0.031) \\
White $-$ Black & $0.087^{*}$ & $0.137^{***}$ \\
 & (0.048) & (0.050) \\
Children (Black / White) & 1,415 / 2,253 & 1,477 / 2,295 \\
\addlinespace
\multicolumn{3}{l}{\textit{Age 35--39}} \\
Pooled (weighted) & $0.454^{***}$ & $0.501^{***}$ \\
 & (0.031) & (0.032) \\
Black & $0.259^{***}$ & $0.313^{***}$ \\
 & (0.044) & (0.048) \\
White & $0.431^{***}$ & $0.464^{***}$ \\
 & (0.037) & (0.035) \\
White $-$ Black & $0.173^{***}$ & $0.151^{**}$ \\
 & (0.058) & (0.060) \\
Children (Black / White) & 1,068 / 1,838 & 1,125 / 1,884 \\
\addlinespace
\multicolumn{3}{l}{\textit{Age 40--44}} \\
Pooled (weighted) & $0.492^{***}$ & $0.530^{***}$ \\
 & (0.037) & (0.034) \\
Black & $0.197^{***}$ & $0.246^{***}$ \\
 & (0.050) & (0.055) \\
White & $0.491^{***}$ & $0.489^{***}$ \\
 & (0.044) & (0.043) \\
White $-$ Black & $0.294^{***}$ & $0.244^{***}$ \\
 & (0.067) & (0.070) \\
Children (Black / White) & 751 / 1,411 & 808 / 1,472 \\
\bottomrule
\end{tabular}}}
\par\smallskip
\begin{minipage}{\textwidth}\footnotesize \textit{Continued on the next table; notes below its last part.}\end{minipage}
\end{table}
\begin{table}[!htbp]
\centering
\addtocounter{table}{-1}
\caption{Log IGE of sons' income by age at measurement (PSID) (continued)}
\resizebox{\ifdim\width>\textwidth\textwidth\else\width\fi}{!}{%
\resizebox*{!}{\ifdim\totalheight>0.85\textheight 0.85\textheight\else\totalheight\fi}{%
\begin{tabular}{lcc}
\toprule
Age band / group & \begin{tabular}[b]{@{}c@{}}ln own earnings\\(child-years)\end{tabular} & \begin{tabular}[b]{@{}c@{}}ln family income\\(child-years)\end{tabular} \\
 & (1) & (2) \\
\midrule
\multicolumn{3}{l}{\textit{Age 45--49}} \\
Pooled (weighted) & $0.484^{***}$ & $0.481^{***}$ \\
 & (0.051) & (0.049) \\
Black & $0.213^{***}$ & $0.215^{***}$ \\
 & (0.064) & (0.072) \\
White & $0.511^{***}$ & $0.473^{***}$ \\
 & (0.060) & (0.061) \\
White $-$ Black & $0.299^{***}$ & $0.258^{***}$ \\
 & (0.088) & (0.095) \\
Children (Black / White) & 496 / 1,001 & 557 / 1,051 \\
\addlinespace
\multicolumn{3}{l}{\textit{Age 50--55}} \\
Pooled (weighted) & $0.498^{***}$ & $0.534^{***}$ \\
 & (0.063) & (0.056) \\
Black & $0.269^{***}$ & $0.290^{***}$ \\
 & (0.074) & (0.089) \\
White & $0.487^{***}$ & $0.478^{***}$ \\
 & (0.073) & (0.070) \\
White $-$ Black & $0.218^{**}$ & $0.187^{*}$ \\
 & (0.103) & (0.114) \\
Children (Black / White) & 351 / 726 & 420 / 780 \\
\addlinespace
\multicolumn{3}{l}{\textit{Age 31--37 (window of Chetty et al.)}} \\
Pooled (weighted) & $0.424^{***}$ & $0.475^{***}$ \\
 & (0.027) & (0.028) \\
Black & $0.292^{***}$ & $0.303^{***}$ \\
 & (0.036) & (0.041) \\
White & $0.392^{***}$ & $0.437^{***}$ \\
 & (0.032) & (0.031) \\
White $-$ Black & $0.099^{**}$ & $0.134^{***}$ \\
 & (0.048) & (0.051) \\
Children (Black / White) & 1,399 / 2,205 & 1,460 / 2,246 \\
\bottomrule
\end{tabular}}}
\par\smallskip
\begin{minipage}{\textwidth}\footnotesize \textit{Notes:} Each cell is the coefficient on log parental family income (mean over the child's ages 13--19) from a separate regression of the outcome in the column header, within the age band of the panel. Child-year columns: child-years in the band, controls for a quartic in the parent's age and single-year age dummies; band-mean columns (daughters' table; for sons in \texttt{output/ige\_by\_age.md}): one observation per child, log of the band mean ($\geq 2$ years), quartic in the parent's age and the mean age. Black and White: unweighted; Pooled: all races, weighted by the core longitudinal weight. White $-$ Black is the difference of the two race-specific regressions, estimated jointly. Standard errors clustered by 1968 family in parentheses. Race is the race of the parental household head. $^{***}$ $p<0.01$, $^{**}$ $p<0.05$, $^{*}$ $p<0.1$. Source: \texttt{26.IGEByAge.py} (\texttt{output/ige\_by\_age.csv}); table built by \texttt{28.PaperBTables.py}.\end{minipage}
\end{table}


\subsection{The racial comparison along the profile}

The general rise hides two profiles. The elasticity of White sons' earnings climbs from 0.221 at 25--29 to 0.431 at
35--39 and 0.491 at 40--44; that of Black sons is 0.262 at 25--29 and between 0.20 and 0.28 at every later age. The
White--Black gap is $-0.041$ (0.045) at 25--29, 0.087$^{*}$ at 30--34, 0.173$^{***}$ at 35--39, 0.294$^{***}$ at
40--44 and 0.299$^{***}$ at 45--49. Its opening, the mean gap from 35 minus the mean gap below 35, is 0.232$^{***}$
(0.052) on bands 25--49, the definition used in the rest of the paper; the White elasticity rises by 0.183$^{***}$ and
the Black one does not move ($-0.049$, s.e.\ 0.039). For family income the profile is the same and the opening smaller,
0.157$^{***}$ (0.055).\footnote{On bands 25--55 (Table~\ref{tab:agetests}) the openings are 0.223$^{***}$ (0.055) and
0.150$^{**}$ (0.060). Holding the children fixed in a panel observed in every band of 25--49, the trend of the
family-income gap is 0.101$^{**}$ (0.046) per ten years.} At 50--55 the earnings gap is 0.218$^{**}$ (0.103): there is no
evidence that it closes, and little power to tell.

\begin{table}[!htbp]
\centering
\caption{Does the log IGE of sons, and the racial gap in it, change with age? (PSID)}
\label{tab:agetests}
\resizebox{\ifdim\width>\textwidth\textwidth\else\width\fi}{!}{%
\resizebox*{!}{\ifdim\totalheight>0.85\textheight 0.85\textheight\else\totalheight\fi}{%
\begin{tabular}{lcc}
\toprule
Panel / statistic & \begin{tabular}[b]{@{}c@{}}ln own earnings\end{tabular} & \begin{tabular}[b]{@{}c@{}}ln family income\end{tabular} \\
 & (1) & (2) \\
\midrule
\multicolumn{3}{l}{\textit{A. Pooled (weighted), all sons}} \\
Equal in all bands: Wald $\chi^2$ [p-value] & $\chi^2(5)=45.27^{***}$ & $\chi^2(5)=62.97^{***}$ \\
 & [$<$0.001] & [$<$0.001] \\
Linear trend per 10 years of age & $0.078^{***}$ & $0.072^{***}$ \\
 & (0.022) & (0.020) \\
Mean over bands below 35 & $0.340^{***}$ & $0.376^{***}$ \\
 & (0.022) & (0.022) \\
Mean over bands 35 and above & $0.482^{***}$ & $0.511^{***}$ \\
 & (0.037) & (0.036) \\
Opening: mean 35 and above minus mean below 35 & $0.142^{***}$ & $0.135^{***}$ \\
 & (0.032) & (0.030) \\
\addlinespace
\multicolumn{3}{l}{\textit{B. Black sons}} \\
Equal in all bands: Wald $\chi^2$ [p-value] & $\chi^2(5)=3.75$ & $\chi^2(5)=5.03$ \\
 & [0.586] & [0.412] \\
Linear trend per 10 years of age & $-0.013$ & $-0.008$ \\
 & (0.029) & (0.033) \\
Mean over bands below 35 & $0.272^{***}$ & $0.274^{***}$ \\
 & (0.032) & (0.033) \\
Mean over bands 35 and above & $0.234^{***}$ & $0.266^{***}$ \\
 & (0.046) & (0.055) \\
Opening: mean 35 and above minus mean below 35 & $-0.038$ & $-0.008$ \\
 & (0.041) & (0.047) \\
\addlinespace
\multicolumn{3}{l}{\textit{C. White sons}} \\
Equal in all bands: Wald $\chi^2$ [p-value] & $\chi^2(5)=50.23^{***}$ & $\chi^2(5)=61.80^{***}$ \\
 & [$<$0.001] & [$<$0.001] \\
Linear trend per 10 years of age & $0.102^{***}$ & $0.074^{***}$ \\
 & (0.026) & (0.026) \\
Mean over bands below 35 & $0.295^{***}$ & $0.335^{***}$ \\
 & (0.027) & (0.026) \\
Mean over bands 35 and above & $0.480^{***}$ & $0.476^{***}$ \\
 & (0.043) & (0.045) \\
Opening: mean 35 and above minus mean below 35 & $0.185^{***}$ & $0.141^{***}$ \\
 & (0.036) & (0.037) \\
\addlinespace
\multicolumn{3}{l}{\textit{D. White $-$ Black}} \\
Equal in all bands: Wald $\chi^2$ [p-value] & $\chi^2(5)=26.79^{***}$ & $\chi^2(5)=23.05^{***}$ \\
 & [$<$0.001] & [$<$0.001] \\
Linear trend per 10 years of age & $0.115^{***}$ & $0.082^{**}$ \\
 & (0.038) & (0.042) \\
Mean over bands below 35 & $0.023$ & $0.060$ \\
 & (0.041) & (0.042) \\
Mean over bands 35 and above & $0.246^{***}$ & $0.210^{***}$ \\
 & (0.063) & (0.071) \\
Opening: mean 35 and above minus mean below 35 & $0.223^{***}$ & $0.150^{**}$ \\
 & (0.055) & (0.060) \\
\bottomrule
\end{tabular}}}
\par\smallskip
\begin{minipage}{\textwidth}\footnotesize \textit{Notes:} Tests on the band-specific log IGEs of Table~\ref{tab:ige_sons} (child-year design; bands 25--29 to 50--55, the 31--37 window excluded), estimated jointly across bands in one stacked, fully interacted regression with standard errors clustered by 1968 family. Wald test of equality across bands (p-value in brackets); linear trend in the band midpoint per 10 years; mean of the band coefficients from 35 minus their mean below 35 (the opening). Bands 25--55 here; the robustness checks of Table~\ref{tab:gaprobust} use bands 25--49 (opening 0.232$^{***}$ for earnings, 0.157$^{***}$ for family income). Race-specific weights and balanced panels: Table~\ref{tab:gaprobust} and Appendix~\ref{app:checks}. $^{***}$ $p<0.01$, $^{**}$ $p<0.05$, $^{*}$ $p<0.1$. Source: \texttt{IGE/26.IGEByAge.py} (\texttt{output/ige\_by\_age.csv}); table built by \texttt{IGE/28.PaperBTables.py}.\end{minipage}
\end{table}


\subsection{How precise is the racial divergence?}

\paragraph{Four data sets.} The PSID estimates rest on 500 to 1,700 Black sons per band. Table~\ref{tab:threedata}
compares four public data sets on ages 25--44. On each data set's own parental measure the earnings gap opens by
0.211$^{***}$ in the PSID and by 0.042 and 0.058 in the NLSY79 and NLSY97; these differ (Cochran $p=0.015$) and are not
pooled. Much of the difference is how parental income is measured: the PSID averages up to seven years, the NLSY97 has
one year and the NLSY79 a three-year mean that is close to one year. On the same PSID children, replacing the 13--19
mean by one year lowers the opening of the family-income gap from 0.162 to 0.093 (difference $-0.069^{**}$), and
instrumenting the single year with parental income at later ages restores it (0.177$^{***}$).\footnote{\texttt{output/age\_profile\_robustness.md},
Section 2. Code: \texttt{33.AgeProfileRobustness.py}.} With one year of parental income in every data set, the
earnings gap opens by 0.151$^{***}$ in the PSID, 0.088$^{**}$ in the NLSY79 and 0.058 (0.040) in the NLSY97; dividing each
race's elasticity by its own data set's reliability ratio gives 0.260$^{***}$, 0.160$^{**}$ and 0.131 (0.086), which
cannot be distinguished ($p=0.47$) and pool to 0.191$^{***}$ (0.046). For family income the corrected openings pool to
0.151$^{***}$ (0.046).

The children of the NLSY79 mothers do not show the opening. Their earnings gap is zero at 25--29 and 30--34, as in the
other data sets ($-0.019$ and 0.025 with the mother's 13--19 mean), and it is $-0.120$ (0.123) at 35--39 and $-0.080$
(0.165) at 40--44, an opening of $-0.103$ (0.106). These later bands rest on 229 Black and 269 White sons at 35--39 and
120 and 93 at 40--44, born by the early 1980s to mothers aged 23 or younger, and the multi-year parental measure does
not raise the opening relative to one year on the same children (difference 0.005, s.e.\ 0.068).\footnote{Code:
\texttt{46.AgeProfileRound2.py} (\texttt{output/age\_profile\_round2.md}, Section A8).} With one year of parental
income in all four data sets, the opening pools to 0.086$^{***}$ (0.024), and equality across the four is rejected only
at 10 percent ($p=0.079$). That is the precision of the racial result: about 0.09 log points with one year of parental
income in four public data sets, about 0.19 corrected for its unreliability in the three where the ratio is known, and
0.15--0.23 on the PSID's multi-year average.

\begin{table}[!htbp]
\centering
\caption{The opening of the racial gap in four public data sets (sons, ages 25--44)}
\label{tab:threedata}
\resizebox{\ifdim\width>\textwidth\textwidth\else\width\fi}{!}{%
\resizebox*{!}{\ifdim\totalheight>0.85\textheight 0.85\textheight\else\totalheight\fi}{%
\begin{tabular}{lccccc}
\toprule
 & \begin{tabular}[b]{@{}c@{}}PSID\end{tabular} & \begin{tabular}[b]{@{}c@{}}NLSY79\end{tabular} & \begin{tabular}[b]{@{}c@{}}NLSY97\end{tabular} & \begin{tabular}[b]{@{}c@{}}NLSY79\\children\end{tabular} & \begin{tabular}[b]{@{}c@{}}Pooled\end{tabular} \\
 & (1) & (2) & (3) & (4) & (5) \\
\midrule
\multicolumn{6}{l}{\textit{Own earnings}} \\
One year of parental income & $0.151^{***}$ & $0.088^{**}$ & $0.058$ & $-0.098$ & $0.086^{***}$ \\
 & (0.042) & (0.045) & (0.040) & (0.094) & (0.024) \\
One year, corrected by each data set's reliability ratio & $0.260^{***}$ & $0.160^{**}$ & $0.131$ &  & $0.191^{***}$ \\
 & (0.073) & (0.081) & (0.086) &  & (0.046) \\
Each data set's main parental measure & $0.211^{***}$ & $0.042$ & $0.058$ & $-0.103$ &  \\
 & (0.048) & (0.042) & (0.040) & (0.106) &  \\
Equality across data sets &  &  &  &  & one year: 4 data sets, $Q$ p=0.079; corrected: $Q$=1.53, p=0.466 \\
\addlinespace
\multicolumn{6}{l}{\textit{Family income}} \\
One year of parental income & $0.088^{**}$ & $0.072$ & $0.079^{*}$ &  & $0.080^{***}$ \\
 & (0.041) & (0.045) & (0.043) &  & (0.025) \\
One year, corrected by each data set's reliability ratio & $0.149^{**}$ & $0.130$ & $0.179^{*}$ &  & $0.151^{***}$ \\
 & (0.072) & (0.082) & (0.092) &  & (0.046) \\
Each data set's main parental measure & $0.137^{***}$ & $0.064$ & $0.079^{*}$ &  &  \\
 & (0.048) & (0.045) & (0.043) &  &  \\
Equality across data sets &  &  &  &  & one year: 3 data sets, $Q$ p=0.965; corrected: $Q$=0.16, p=0.923 \\
\bottomrule
\end{tabular}}}
\par\smallskip
\begin{minipage}{\textwidth}\footnotesize \textit{Notes:} Opening of the White $-$ Black gap in the log IGE: mean gap at 35--39 and 40--44 minus mean gap at 25--29 and 30--34, child-year design, clustered by family (CYA: mother). One year of parental income: PSID the wave at child age 12--16 nearest 14, NLSY79 1978 income, NLSY97 1996 income, NLSY79 children: the mother's family income in the year nearest the child's age 14. Corrected: each race's IGE divided by that data set's reliability ratio of one year of family income (PSID White 0.585, Black 0.565; NLSY79 0.552, 0.556; NLSY97 0.430, 0.492), with the ratios' sampling error propagated. Main measures: PSID mean at 13--19, NLSY79 1978--80 mean, NLSY97 1996, NLSY79 children the mother's mean family income at the child's 13--19; main measures are not pooled because the four differ in reliability. NLSY79 children: wage and salary income, children of NLSY79 mothers, few observed at 35--44 (93 White and 120 Black sons at 40--44, born when their mothers were 23 or younger). Pooled: inverse-variance weighted. $^{***}$ $p<0.01$, $^{**}$ $p<0.05$, $^{*}$ $p<0.1$. Source: \texttt{41.AgeProfileMeasurement.py} (\texttt{output/age\_profile\_measurement.csv}), \texttt{46.AgeProfileRound2.py} and \texttt{33.AgeProfileRobustness.py}.\end{minipage}
\end{table}


\paragraph{Sample, attrition, period, experience and metric.} Table~\ref{tab:gaprobust} collects the PSID checks on the
opening, bands 25--49, earnings first. On the SRC sample alone, which has 217 Black sons, the opening is 0.279$^{***}$;
with each race weighted, 0.227$^{***}$; with inverse-probability weights for staying in the panel, 0.201$^{***}$
(Section~\ref{sec:data}).

Age, period and cohort are collinear within a cohort, and the log gap could reflect rising inequality over calendar
time \citep{bayer2018divergent}. Three pieces of evidence separate age from period. The opening of the earnings gap is
between 0.17 and 0.28, and significant, in each of the cohorts born 1951--64, 1965--79 and 1980--89, which are observed
in different decades.\footnote{\texttt{output/age\_profile\_robustness.md}, Section 1.} Income-year fixed effects in every
band and race cell, a period-varying elasticity evaluated in 1995 and a fixed calendar window (income years 1990--2010,
in which the bands hold different cohorts) give 0.243$^{***}$, 0.235$^{***}$ and 0.216$^{***}$. And within fixed cohorts
the dispersion of White sons' log earnings rises with age---from 0.67 at 25--29 to 0.83 at 45--49 for sons born 1951--64,
and from 0.75 to 0.96 for sons born 1965--79---while that of Black sons does not (0.93 to 0.79, and 0.84 to 0.84); the
White--Black difference in the rise is 0.149$^{***}$ and 0.106$^{**}$.\footnote{Within-race standard deviation of log
earnings net of race-specific single-age means, bootstrap standard errors over families;
\texttt{output/age\_profile\_measurement.md}, Section B(iv). Code: \texttt{41.AgeProfileMeasurement.py}.} The
family-income opening is less robust to the calendar: 0.174$^{***}$ with year effects, 0.168 ($p=0.106$) with a
period-varying elasticity and 0.084 ($p=0.32$) in the fixed window.

White sons of richer parents enter the labor market later: at 25--29, sons of top-tercile White parents have 4.8 years
of potential experience and other sons 6.0--7.0. Measured on an experience clock, the gap is $-0.121$ (0.101) at 0--4
years, 0.043 (0.044) at 5--9, 0.118$^{**}$ at 10--14, 0.220$^{***}$ at 15--19 and 0.272$^{***}$ at 20--24 (Appendix
Table~\ref{tab:experience}), and it opens by 0.243$^{***}$ between the first and the second decade of the career; dropping
the years at 25--26 leaves 0.212$^{***}$. The zero gap of the late twenties is not a difference in years of experience:
the gap is zero for the first ten years of the career and opens in the second.

Log elasticities drop years without earnings and are sensitive to the bottom of the distribution
\citep{mitnik2015new, mitnik2020intergenerational}. The dollar-weighted elasticity, estimated on the earnings level with
zeros kept, opens by 0.257$^{***}$. In the measurement literature's own terms \citep{haider2006life,
nybom2016heterogeneous}, the coefficient of current on lifetime earnings rises with age from 0.69 at 25--29 to 1.22 at
45--55 for White sons and from 0.85 to 0.99--1.02 for Black sons, a difference in the rise of 0.211$^{**}$ (0.086), and
parental income adds nothing to current earnings once lifetime earnings are held fixed, in either race (Appendix
Table~\ref{tab:lambda}).\footnote{Sons with at least three valid years in each of 25--34, 35--44 and 45--55 (267 Black
and 676 White sons); a long panel is the definition of this object, not a test of the age profile. Within parental
terciles the racial difference in the rise is 0.04--0.17 and not significant, which is consistent both with low power
and with the position reading below.}

\begin{table}[!htbp]
\centering
\caption{How robust is the opening of the racial gap? (PSID, sons)}
\label{tab:gaprobust}
\resizebox{\ifdim\width>\textwidth\textwidth\else\width\fi}{!}{%
\resizebox*{!}{\ifdim\totalheight>0.85\textheight 0.85\textheight\else\totalheight\fi}{%
\begin{tabular}{lcc}
\toprule
 & \begin{tabular}[b]{@{}c@{}}Own earnings\end{tabular} & \begin{tabular}[b]{@{}c@{}}Family income\end{tabular} \\
 & (1) & (2) \\
\midrule
\multicolumn{3}{l}{\textit{Baseline}} \\
Baseline: child-year design, bands 25--49 & $0.232^{***}$ & $0.157^{***}$ \\
 & (0.052) & (0.055) \\
\addlinespace
\multicolumn{3}{l}{\textit{Sample and attrition}} \\
SRC sample only (no SEO) & $0.279^{***}$ & $0.313^{***}$ \\
 & (0.095) & (0.091) \\
Each race weighted & $0.227^{***}$ & $0.184^{**}$ \\
 & (0.069) & (0.072) \\
Sons observed at 25--29 & $0.213^{***}$ & $0.143^{**}$ \\
 & (0.057) & (0.061) \\
\quad + inverse-probability weights for staying & $0.201^{***}$ & $0.129^{**}$ \\
 & (0.063) & (0.065) \\
\addlinespace
\multicolumn{3}{l}{\textit{Age, period and cohort}} \\
Cohort-varying IGE (1965 cohort) & $0.240^{***}$ & $0.176^{***}$ \\
 & (0.056) & (0.058) \\
Income-year fixed effects & $0.243^{***}$ & $0.174^{***}$ \\
 & (0.052) & (0.056) \\
Period-varying IGE (1995) & $0.235^{***}$ & $0.168$ \\
 & (0.087) & (0.104) \\
Income years 1990--2010 only & $0.216^{***}$ & $0.084$ \\
 & (0.076) & (0.084) \\
\addlinespace
\multicolumn{3}{l}{\textit{Experience clock}} \\
Potential-experience bands ($\geq 10$ minus $<10$ years) & $0.243^{***}$ & $0.227^{***}$ \\
 & (0.067) & (0.065) \\
Ages 27--49 (years at 25--26 dropped) & $0.212^{***}$ & $0.129^{**}$ \\
 & (0.052) & (0.056) \\
\addlinespace
\multicolumn{3}{l}{\textit{Metric}} \\
Dollar-weighted IGE (PPML, zeros kept) & $0.257^{***}$ & $0.145^{***}$ \\
 & (0.061) & (0.052) \\
\addlinespace
\multicolumn{3}{l}{\textit{Position in the parental distribution}} \\
Parents below the pooled median & $0.130$ & $0.023$ \\
 & (0.087) & (0.086) \\
Parents in the bottom pooled tercile & $0.232^{*}$ & $0.092$ \\
 & (0.127) & (0.121) \\
White sons re-weighted to Black parents & $0.188^{*}$ & $0.079$ \\
 & (0.100) & (0.087) \\
\addlinespace
\multicolumn{3}{l}{\textit{Ranks of family income (column 1: sons; column 2: sons and daughters)}} \\
Rank--rank slope, national ranks & $0.037$ & $0.080$ \\
 & (0.077) & (0.049) \\
Rank--rank slope, within-race ranks & $0.015$ & $0.093^{**}$ \\
 & (0.057) & (0.043) \\
\bottomrule
\end{tabular}}}
\par\smallskip
\begin{minipage}{\textwidth}\footnotesize \textit{Notes:} Each cell is the change of the White $-$ Black gap in the IGE from the age bands below 35 to the bands from 35 (mean gap over 35--39, 40--44, 45--49 minus mean gap over 25--29, 30--34), with its clustered standard error (1968 family), from stacked band $\times$ race regressions of the child-year design (log outcome on log parental income, quartic in the parent's age, single-year age dummies). Attrition: sons with family income observed at 25--29; weights are the inverse of the predicted probability of being observed in each later band (logit on log parental income, income at 25--29, both interacted with race, SRC, parent's age, five-year cohorts). Period-varying: log parental income $\times$ (year $-$ 1995)/10 in every cell. Experience clock: bands of potential experience (age $-$ years of schooling $-$ 6; 0--4 to 20--24) with single-year experience dummies, opening = mean gap at 10 or more years minus mean gap below 10; or the age profile without the years at 25--26. PPML: Poisson pseudo-maximum likelihood on the income level. Position: pooled quantiles of parental income over Black and White sons; re-weighting by the ratio of Black to White shares in 20 parental-income bins, band by band. Ranks: rank--rank slope of band-mean family income (column 2: sons and daughters). $^{***}$ $p<0.01$, $^{**}$ $p<0.05$, $^{*}$ $p<0.1$. Source: \texttt{41.AgeProfileMeasurement.py} (\texttt{output/age\_profile\_measurement.csv}), \texttt{46.AgeProfileRound2.py} and \texttt{33.AgeProfileRobustness.py}.\end{minipage}
\end{table}


\subsection{Logs, ranks and position in the distribution}

\paragraph{Ranks.} The PSID reproduces \citet{chetty2020race} at 31--37: rank--rank slopes of 0.370 for Black and
0.385 for White children, and a gap of 12.9$^{***}$ (1.3) percentiles at the 25th parental percentile. For sons' family income the rank
evidence does not change with age: a gap in rank--rank slopes cannot be rejected in any band, and its opening is 0.037
(0.077) in national ranks and 0.015 (0.057) when sons and parents are ranked within race; the gap in expected rank at
the 25th parental percentile is 8--11 percentiles in every band (Appendix Table~\ref{tab:rankage}). In sons' own
earnings the rank slopes do diverge after 40 (0.150$^{*}$ at 40--44, column (4) of the same table), and in hourly
wages on the sample of the model by 0.187$^{***}$ (0.062, $p = 0.003$) at 45--49, against 0.052 (0.036) at 25--29
(Table~\ref{tab:modeluntargeted}): ranks of family income do
not register the racial divergence of sons' own pay. For sons and
daughters together the within-race rank slope gap does open, by 0.093$^{**}$ (0.043; national ranks 0.080, s.e.\ 0.049),
so ranks are not immune to the age profile; they register a smaller part of it.

\paragraph{Why logs and ranks differ.} On the same children and weights the log gap widens by 0.216$^{***}$ (0.062)
between the two halves of the career and the rank gap by 0.113$^{**}$ (0.049); the widening of the log gap is mostly
the dispersion of children's log income relative to their parents' (the gap in the ratio of standard deviations widens
by 0.262$^{***}$), which ranks normalise away.\footnote{Both sexes, family income, \texttt{output/ige\_by\_age.md},
Section 5.4. Code: \texttt{26.IGEByAge.py}.} In Chetty et al.'s language, the log gap is the downward mobility of
Black sons from well-off families: in the PSID 17.8 percent of Black and 5.9 percent of White children of top-quintile
parents end in the bottom quintile.\footnote{Code: \texttt{7.RankMobility.py}.} That downward mobility is known
\citep{hertz2005rags, bhattacharya2011nonparametric, mazumder2014black}; the new fact is its timing.

\paragraph{Race or position?} Black families are concentrated at the bottom of the parental distribution---85 percent
of Black and 46 percent of White sons have parents below the pooled median---where log incomes are dispersed and the
IGE is flatter \citep{bratsberg2007nonlinearities}. Figure~\ref{fig:position} looks at the question nonparametrically.
For own earnings (panel A), White sons' earnings at 40--49 rise more steeply with parental income than at 25--29 across
the whole common range, including the bottom six vigintiles, which hold 54--316 Black and 20--125 White sons each; Black
sons' 40--49 curve is no steeper than their 25--29 curve, and above the sixth vigintile it rests on 15--33 Black sons
per point. The regressions are less sharp than the figure. Among sons of parents in the bottom tercile, below the median
or with White sons re-weighted to the parental incomes of Black sons, the earnings opening is 0.13--0.23 and the
family-income opening 0.02--0.09, with standard errors of 0.09--0.13 (panels C and D and Table~\ref{tab:gaprobust}). Both
are smaller than in the full sample, and the data do not establish whether the difference between the two outcomes, or
between these subsamples and the full sample, is position or race. The model of Section~\ref{sec:model} lets access to
good jobs depend nonlinearly on parental income, so that position can do this work before race is asked to.

\begin{figure}[!htbp]
\centering
\includegraphics[width=\textwidth]{paperB_position.pdf}
\caption{Race or position in the parental distribution? (PSID, sons)}\label{fig:position}
\begin{minipage}{\textwidth}\footnotesize \textit{Notes:} Panels A and B: mean log outcome (net of race-specific
single-age means) by vigintile of the pooled parental-income distribution of Black and White sons, at 25--29 (dashed)
and 40--49 (solid); vigintiles with at least 20 sons of the race. Panels C and D: White $-$ Black gap in the log IGE by
age band, all sons, sons of parents below the pooled median, and White sons re-weighted to the Black distribution of
parental income in each band (20 bins); 95\% confidence intervals clustered by 1968 family. Code:
\texttt{41.AgeProfileMeasurement.py} and \texttt{44.PaperBReviewTables.py}.\end{minipage}
\end{figure}

\paragraph{Daughters.} Daughters' family income shows the same divergence (gap trend 0.162$^{***}$, s.e.\ 0.032),
their own earnings do not (Appendix Table~\ref{tab:ige_daughters}). The daughters' divergence runs through their
husbands (Section~\ref{sec:contains}), and it rests on the low-income oversample: on the SRC sample alone it is 0.014
(0.106).\footnote{\texttt{output/age\_profile\_robustness.md}, Section 3. Code: \texttt{33.AgeProfileRobustness.py}.}

The age profile raises the question the rest of the paper addresses: what grows with parental income over White sons'
thirties and not over Black sons'?
